\documentclass[11pt]{article}
\usepackage[utf8]{inputenc}
\usepackage{amsmath,amssymb}
\usepackage[margin=1in]{geometry}
\usepackage{hyperref}
\emergencystretch=3em

\title{Continuous test functions from monomial moments}
\author{Claude, with Daniel Murfet}
\date{2026-09-07}

\begin{document}
\maketitle
\sloppy

\begin{abstract}
Second step of the continuous-amplitude milestone (Astra~\#16, route~C), independent of the
Laplace-specific constants. On the closed unit cube $K=[0,1]^d$: every continuous $\eta$ is uniformly
approximated by a multivariate polynomial (\texttt{exists\_mvPolynomial\_near}, Stone--Weierstrass for
the subalgebra generated by the coordinate functions, whose elements are exactly polynomial functions),
and nonnegative weighted functionals $T_N(\eta)=\int_S\eta(\varphi_N x)\,w_N(x)\,dx$ on $S\subseteq K$
with maps $\varphi_N,\varphi_L$ into $K$ converge against every continuous test function as soon as
they converge on monomials (\texttt{tendsto\_integral\_of\_monomials}). Zero \texttt{sorry}/\texttt{axiom}.
\end{abstract}

\section{The things formalised}
{\small
\begin{verbatim}
def closedCube (d : ℕ) : Set (Fin d → ℝ) := Set.pi univ fun _ => Icc 0 1
def coordAlgebra (d : ℕ) : Subalgebra ℝ C(Fin d → ℝ, ℝ) :=
  Algebra.adjoin ℝ (Set.range fun i => (ContinuousMap.eval i : C(Fin d → ℝ, ℝ)))
theorem coordAlgebra_separatesPoints (d) : (coordAlgebra d).SeparatesPoints
theorem mem_coordAlgebra_eval (hg : g ∈ coordAlgebra d) :
    ∃ p : MvPolynomial (Fin d) ℝ, ∀ x, g x = MvPolynomial.eval x p
theorem exists_mvPolynomial_near (d) (η) (hη : Continuous η) (ε) (hε : 0 < ε) :
    ∃ p : MvPolynomial (Fin d) ℝ, ∀ x ∈ closedCube d, |η x - MvPolynomial.eval x p| < ε
theorem tendsto_of_approx (f : ℝ → ℝ) (a : ℝ) (h : ∀ ε > 0, ∃ g b, Tendsto g atTop (𝓝 b)
    ∧ (∀ᶠ N in atTop, |f N - g N| ≤ ε) ∧ |a - b| ≤ ε) : Tendsto f atTop (𝓝 a)
theorem integrableOn_comp_mul_of_continuous ... : IntegrableOn (fun x => f (φ x) * w x) S
theorem integral_eval_mul_eq_sum ... : ∫ x in S, MvPolynomial.eval (φ x) p * w x
    = ∑ γ ∈ p.support, MvPolynomial.coeff γ p * ∫ x in S, (∏ i, φ x i ^ γ i) * w x
theorem tendsto_integral_of_monomials (d S hS φ φL hφ hφL hφm hφLm w wL hw hwL hwnn hwLnn)
    (hmono : ∀ γ : Fin d → ℕ, Tendsto (fun N => ∫ x in S, (∏ i, φ N x i ^ γ i) * w N x)
      atTop (𝓝 (∫ x in S, (∏ i, φL x i ^ γ i) * wL x))) (η) (hη : Continuous η) :
    Tendsto (fun N => ∫ x in S, η (φ N x) * w N x) atTop (𝓝 (∫ x in S, η (φL x) * wL x))
\end{verbatim}
}

\section{Mathematics}
The constant monomial gives $T_N(1)\to T_L(1)=:M$, hence an eventual mass bound $T_N(1)\le M+1$.
For a polynomial $p=\sum_\gamma c_\gamma u^\gamma$, finite linearity of the integral reduces
$T_N(p)\to T_L(p)$ to the monomial hypothesis. For continuous $\eta$ and $\varepsilon>0$, choose $p$
with $\sup_K|\eta-p|<\delta:=\varepsilon/(M+2)$; positivity of the weights gives
$|T_N(\eta)-T_N(p)|\le\delta\,T_N(1)\le\delta(M+1)\le\varepsilon$ eventually and
$|T_L(\eta)-T_L(p)|\le\delta M\le\varepsilon$; the $\varepsilon/3$ lemma finishes.

\section{Paper-facing meaning}
This is the analytic device that upgrades the shifted monomial moments of unit~184 to continuous
amplitudes: the limiting functional may be supported on the face $u_J=0$ ($\varphi_L$ the projection
killing the minimal coordinates) while the finite-$N$ functionals use $\varphi_N=\mathrm{id}$.

\section{Lean notes}
\begin{itemize}
\item \texttt{∞} is a token (the \texttt{ENNReal} top), so \texttt{φ∞} is not an identifier; use
\texttt{φL}.
\item Stone--Weierstrass: \texttt{ContinuousMap.exists\_mem\_subalgebra\_near\_}\allowbreak\texttt{continuous\_of\_isCompact\_of\_separatesPoints}
with the adjoin of the coordinate evaluations; \texttt{Algebra.adjoin\_induction} via
\texttt{induction hg using ... with | mem | algebraMap | add | mul} produces the polynomial.
\item \texttt{MvPolynomial.continuous\_eval} takes the polynomial explicitly and lives in
\texttt{Mathlib.Topology.Algebra.MvPolynomial}, which must be imported explicitly.
\item Sums over \texttt{p.support} index by \texttt{Fin d →₀ ℕ}; the monomial hypothesis quantifies over
\texttt{Fin d → ℕ}, so instantiate it at \texttt{fun i => γ i} and annotate the binder.
\end{itemize}

\end{document}
